% \begin{document}
\chapter{Access Control}

\section{Autenticazione e autorizzazione}
\dfn{Autenticazione}{
  Proprieta' di un sistema di garantire che un utente possa essere \textbf{identificato} tramite informazioni in suo possesso. Le identificazioni possono essere di tre tipi (usati da soli o in combinazione):
  \begin{itemize}
    \item \textbf{Quello che si ha} (badge, smart card, \dots) — detto \textit{token}.
    \item \textbf{Quello che si sa} (password, PIN, risposte a domande prestabilite).
    \item \textbf{Quello che si e'} (riconoscimento tramite impronte digitali, retina, viso).
  \end{itemize}
}
Una volta autenticati i soggetti, il problema successivo e' l'\textbf{autorizzazione} o controllo degli accessi.

\section{Controllo degli accessi (Access Control)}
Il controllo degli accessi e' un elemento centrale della sicurezza informatica; i suoi obiettivi sono:
\begin{itemize}
  \item impedire a utenti \textbf{non autorizzati} di accedere alle risorse;
  \item impedire agli utenti legittimi di accedere alle risorse in modo \textbf{non autorizzato};
  \item consentire agli utenti legittimi di accedere alle risorse in modo \textbf{autorizzato}.
\end{itemize}

\dfn{Controllo degli accessi}{
  Insieme di politiche e meccanismi che servono a decidere se un particolare soggetto ha il permesso di eseguire determinate operazioni su determinati oggetti. Gli elementi base sono:
  \begin{itemize}
    \item \textbf{soggetto}: un'entita' in grado di accedere agli oggetti;
    \item \textbf{oggetto}: una risorsa a cui e' necessario controllare l'accesso;
    \item \textbf{diritto di accesso}: descrive il modo in cui un soggetto puo' accedere a un oggetto.
  \end{itemize}
}
La \textbf{politica} di controllo degli accessi stabilisce quali tipi di accesso sono consentiti, in quali circostanze e da chi.

\subsection{Tipi di politiche}
\begin{enumerate}
  \item \textbf{DAC} — Discretionary Access Control (Controllo discrezionale)
  \item \textbf{MAC} — Mandatory Access Control (Controllo obbligatorio)
  \item \textbf{RBAC} — Role-Based Access Control (Controllo basato sui ruoli)
\end{enumerate}

\section{DAC — Discretionary Access Control}
\dfn{DAC}{
  L'accesso e' basato sull'\textbf{identita' dei soggetti} e su \textbf{regole di accesso} che stabiliscono cosa i soggetti possono (o non possono) fare su quali oggetti. E' \textit{discrezionale} perche' i soggetti \textbf{decidono loro} di concedere (o negare) l'accesso ad altri soggetti.
}

\subsection{Matrice di accesso}
Un approccio generale al DAC e' quello di una \textbf{matrice di accesso}: le righe sono i soggetti, le colonne gli oggetti, e ogni cella contiene i diritti di accesso.

\begin{center}
  \footnotesize
  \begin{tabular}{|l|c|c|c|c|}
    \hline
    & \textbf{File 1} & \textbf{File 2} & \textbf{File 3} & \textbf{File 4} \\ \hline
    \textbf{User A} & Own, Read, Write & & Own, Read, Write & \\ \hline
    \textbf{User B} & Read & Own, Read, Write & Write & Read \\ \hline
    \textbf{User C} & Read, Write & Read & & Own, Read, Write \\ \hline
  \end{tabular}
\end{center}

La matrice puo' essere rappresentata in due modi equivalenti:
\begin{itemize}
  \item \textbf{ACL} (Access Control List): si codificano le \textbf{colonne} — per ogni oggetto, la lista dei soggetti autorizzati e dei relativi diritti.
  \item \textbf{Capability}: si codificano le \textbf{righe} — per ogni soggetto, la lista degli oggetti a cui ha accesso e con quali diritti.
\end{itemize}

\subsection{Notazione formale}
Sia $ S $ l'insieme dei \textbf{soggetti}, $ O $ l'insieme degli \textbf{oggetti} (gli oggetti possono essere attivi e agire anche da soggetti), e $ \alpha $ l'insieme dei \textbf{diritti di accesso} che i soggetti hanno sugli oggetti.

\dfn{Protection domain}{
  Un \textbf{dominio di protezione} e' un insieme di oggetti con i relativi diritti di accesso. Formalmente e' un insieme di coppie:
  \[
    \langle \texttt{object},\; \texttt{set\_of\_access\_rights} \rangle
  \]
  I soggetti sono associati a un dominio di protezione in cui operano. L'associazione puo' essere \textbf{statica} (il dominio non cambia) o \textbf{dinamica} (il dominio puo' cambiare a runtime).
}

\ex{Kernel mode vs User mode}{
  "Kernel mode" e "user mode" sono due \textbf{domini di protezione} che controllano l'accesso alla memoria principale. Normalmente i processi operano in user mode; quando eseguono una system call passano in kernel mode e acquisiscono i privilegi necessari. E' un esempio di associazione \textbf{dinamica} tra soggetto e dominio.
}

\subsection{Matrice di accesso — modello formale}
\dfn{Access Control Matrix}{
  Una matrice $ M $ con i \textbf{domini} come righe e gli \textbf{oggetti} come colonne. La cella $ M(i,j) $ contiene l'insieme dei diritti di accesso $ \alpha $ che il dominio $ D_i $ ha sull'oggetto $ O_j $. Quando viene creato un nuovo oggetto si aggiunge una nuova colonna; il contenuto e' deciso dal creatore.
}

\subsubsection*{Il diritto "switch" e il least privilege}
\ex{Dictionary problem}{
  Utente $ A $ (dominio $ D_2 $) modifica $ \mathit{File}_2 $; utente $ B $ (dominio $ D_3 $) modifica $ \mathit{File}_3 $. Entrambi attivano un correttore ortografico basato su $ \mathit{File}_4 $, un dizionario proprietario che non deve poter essere copiato. Se si danno diritti \textit{read} diretti su $ \mathit{File}_4 $, $ A $ e $ B $ potrebbero copiarlo.

  Soluzione: introdurre un dominio $ D_4 $ che puo' \textbf{solo} leggere $ \mathit{File}_4 $, e aggiungere il diritto \textbf{switch}: $ D_2 $ e $ D_3 $ possono passare a $ D_4 $, ma non hanno piu' accesso diretto a $ \mathit{File}_4 $. Lo "switch" non solo cambia dominio ma copia i diritti del dominio sorgente in quello destinazione; poiche' piu' utenti possono passare allo stesso $ D_4 $, si creano istanze logicamente distinte. Questo meccanismo implementa il \textbf{principio di least privilege}.
}

\subsubsection*{Implementazione della matrice}
\begin{center}
  \footnotesize
  \begin{tabular}{|l|p{4.5cm}|p{4.5cm}|}
    \hline
    \textbf{Approccio} & \textbf{Vantaggi} & \textbf{Svantaggi} \\ \hline
    Tabella globale & Semplice da implementare & Tabella enorme; difficile da mantenere in sistemi dinamici \\ \hline
    ACL (per colonna) & Revoca facile (si aggiorna la lista dell'oggetto); controllo di accesso basato sull'identita' & Efficienza minore per accessi ripetuti \\ \hline
    Capability (per riga) & Accesso efficiente (la capability e' gia' in mano al processo) & Revoca difficile (i diritti sono distribuiti ai processi) \\ \hline
  \end{tabular}
\end{center}

\section{MAC — Mandatory Access Control}
\dfn{MAC}{
  L'accesso e' basato sul confronto delle \textbf{etichette di sicurezza} (\textit{security labels}, che indicano quanto sono sensibili o critici gli oggetti) con le \textbf{autorizzazioni di sicurezza} (\textit{security clearances}) dei soggetti. E' \textit{obbligatorio} perche' le etichette e i nulla osta sono impostati dal \textbf{sistema} e non possono essere modificati dai soggetti.
}
I primi sistemi MAC erano basati sul modello \textbf{Bell-LaPadula} (BLP). Su Linux esistono diverse implementazioni: SeLinux, Tomoyo, SMACK, \dots \\
E' facile pensare a questi sistemi pensando alla segregazione delle applicazioni Android: Telegram non puo' vedere i messaggi di WhatsApp.

\section{RBAC — Role-Based Access Control}
\dfn{RBAC}{
  L'accesso e' basato sui \textbf{ruoli} che i soggetti hanno all'interno del sistema e su \textbf{regole} che stabiliscono quali accessi sono consentiti ai soggetti in determinati ruoli. Agli utenti vengono assegnati ruoli diversi, in modo statico o dinamico, in base alle loro responsabilita'.
}
La matrice di accesso in RBAC associa gli utenti ai ruoli (non direttamente agli oggetti). Ogni voce e' vuota o contrassegnata (l'utente e' o non e' in quel ruolo). Un singolo utente puo' avere piu' ruoli; piu' utenti possono condividere un singolo ruolo.

\section{Principi fondamentali per le politiche di sicurezza}
\begin{enumerate}
  \item \textbf{Open design} (\textit{No security by obscurity}): la sicurezza non dovrebbe dipendere dall'ignoranza dei potenziali avversari riguardo ai meccanismi di protezione, bensi' dalla segretezza di chiavi o password specifiche. Analogo al principio di Kerckhoffs.
  \item \textbf{Economy of mechanism} (\textit{as simple as possible}): semplificare la progettazione e l'implementazione dei meccanismi di sicurezza riduce le possibilita' di errori. Interfacce piu' semplici tra moduli di sicurezza riducono la superficie di attacco.
  \item \textbf{Fail-safe defaults}: a meno che a un soggetto non venga concesso esplicitamente l'accesso a un oggetto, gli dovrebbe essere \textbf{negato}. Il default e' il divieto.
  \item \textbf{Complete mediation} (reference monitor): non e' possibile accedere direttamente agli oggetti; \textbf{tutti} gli accessi devono passare per il meccanismo di controllo.
  \item \textbf{Least privilege}: i soggetti dispongono dei \textbf{privilegi di accesso minimi} necessari per effettuare le operazioni necessarie per quella fase di esecuzione.
\end{enumerate}

\subsection{Least Privilege in dettaglio}
Il principio di least privilege:
\begin{itemize}
  \item Limita i danni che possono derivare da un incidente o da un errore.
  \item Limita il numero di programmi privilegiati.
  \item Aiuta nel debug e aumenta la sicurezza.
  \item Consente l'isolamento dei sotto-sistemi critici.
\end{itemize}
Il privilegio minimo e' imposto attraverso un reference monitor che implementa una \textbf{completa mediazione}: ogni accesso a ogni oggetto viene controllato.

In RBAC: ogni ruolo dovrebbe disporre del set minimo di diritti di accesso necessari per quel ruolo. A un utente viene assegnato il ruolo che gli consente di svolgere solo le attivita' necessarie; piu' utenti con lo stesso ruolo godono dello stesso set minimo.

\section{Sicurezza dei sistemi}
La sicurezza dei sistemi differisce da quella delle reti per tre aspetti:
\begin{itemize}
  \item Possibilita' di avere un ente certificato nel mezzo (il Sistema Operativo).
  \item I sistemi possono essere singolo o multi-utente.
  \item Possibilita' di divisione dei privilegi.
\end{itemize}
E' la superficie su cui agiscono i \textbf{malware locali}.

\dfn{Privilege escalation}{
  La possibilita' di ottenere piu' privilegi sul sistema sfruttando una falla, un errore di configurazione di un software applicativo o di un sistema operativo. Esempi: installare programmi, leggere informazioni private o modificare configurazioni del sistema.
}
\dfn{Arbitrary Code Execution (ACE)}{
  Il primo passo per il privilege escalation e' l'esecuzione arbitraria di codice. Puo' essere perpetrata:
  \begin{itemize}
    \item installando software;
    \item facendo girare software;
    \item sfruttando vulnerabilita' in grado di dirottare il funzionamento dei programmi.
  \end{itemize}
}

\section{UNIX e il filesystem}
\nt{In UNIX (piu' o meno!) \textbf{tutto e' un file}. Avendo controllo su un file si puo' ottenere controllo su quello che rappresenta (es.\ \texttt{/dev/snd/*} per l'audio).}

\subsection{Struttura del filesystem Linux}
Il filesystem si compone di una radice (\texttt{/}) dalla quale il sistema si dirama come un albero:

\begin{center}
  \footnotesize
  \begin{tabular}{|l|p{9cm}|}
    \hline
    \texttt{/} & The root filesystem \\ \hline
    \texttt{/proc} & (pseudo) Process virtual infrastructure (es.\ \texttt{/proc/1/cmdline}). \\ \hline
    \texttt{/sys} & (pseudo) System virtual infrastructure (es.\ \texttt{/sys/class/leds}). \\ \hline
    \texttt{/dev} & (pseudo) Device drivers file interface (es.\ \texttt{/dev/sda}). \\ \hline
    \texttt{/run} & (pseudo) Contains run-time information. \\ \hline
    \texttt{/etc} & Configuration directory. \\ \hline
    \texttt{/tmp} & Temporary directory, volatile, usually mounted in RAM. \\ \hline
    \texttt{/root} & Root's home. \\ \hline
    \texttt{/bin} & Binaries. \\ \hline
    \texttt{/lib} & Libraries. \\ \hline
    \texttt{/sbin} & System Binaries. \\ \hline
    \texttt{/usr} & User binaries. \\ \hline
    \texttt{/var} & Log, cache, and structured personal files (es.\ mailboxes). \\ \hline
    \texttt{/home} & Users' directories. \\ \hline
    \texttt{/mnt} & External mountpoint. \\ \hline
    \texttt{/opt} & Optional softwares. \\ \hline
  \end{tabular}
\end{center}

\subsection{Il file \texttt{/etc/passwd}}
Il file \texttt{/etc/passwd} tiene traccia di ogni utente registrato che ha accesso al sistema. E' separato da due punti (\texttt{:}) e contiene, in ordine: nome utente, password crittografata, UID, GID, nome completo (GECOS), directory home, shell di login.

Esempio:
\begin{verbatim}
root:x:0:0:root:/root:/bin/bash
daemon:x:1:1:daemon:/usr/sbin:/usr/sbin/nologin
\end{verbatim}

I permessi di \texttt{/etc/passwd} sono \texttt{-rw-r--r--} (leggibile da tutti, scrivibile solo da root).

\subsection{Il file \texttt{/etc/group}}
Analogamente, \texttt{/etc/group} associa ad ogni gruppo un GID:
\begin{verbatim}
root:x:0:
daemon:x:1:
\end{verbatim}

\section{ACL in Linux}
\dfn{Access Control List (ACL)}{
  Il sistema mantiene le informazioni di accesso ai file tramite una tabella soggetto/oggetto. Ogni riga e' un soggetto; ogni colonna e' un oggetto; le celle contengono i diritti (\texttt{R}, \texttt{W}, \texttt{X}).
}
Unix/Linux (senza le estensioni POSIX ACL) dichiara \textbf{3 soggetti} per ogni file:
\begin{itemize}
  \item \textbf{user}: il proprietario del file;
  \item \textbf{group}: il gruppo proprietario del file;
  \item \textbf{other}: tutti gli altri.
\end{itemize}

Il sistema in modalita' DAC analizza (in questo ordine): UID del processo e proprietario del file; GID del processo e gruppo del file; \dots
\nt{
  Questo puo' generare complicanze contro-intuitive: l'utente appartiene a un gruppo che non puo' leggere il file, ma e' il suo proprietario e il proprietario non puo' leggere il file.
}

\subsection{Il super utente}
Il super utente privilegiato del sistema e' quello che puo' svolgere ogni operazione a partire dalla radice del file system. Si chiama \textbf{root} in Linux/UNIX e \textbf{SYSTEM} in Windows.

\section{Capabilities}
\dfn{Capability}{
  Un \textbf{token} o \textbf{biglietto} in grado di rappresentare un determinato oggetto sul quale si ha l'accesso. La capability specifica gli oggetti e le operazioni autorizzati per un particolare utente. Ogni utente dispone di un numero di ticket e puo' essere autorizzato a prestarli o cederli ad altri.
}
Esempi di capability: i cookie web (identificatori che indicano a quale sessione e' associata una comunicazione), i file aperti.

A seguito di una system call \texttt{open}, viene ritornato un identificativo univoco detto \textbf{file descriptor} che il sistema operativo usa per indicare il file aperto. L'integrita' del ticket deve essere protetta (solitamente dal SO); il ticket deve essere \textbf{irrefutabile}.

Le capability possono essere supportate da sistemi crittografici (cookie) o da segmentazione a livello di sistema operativo (file descriptor).

\nt{
  Unix utilizza un sistema \textbf{misto}: ACL per i file non aperti; Capability (file descriptor, socket, \dots) per i file aperti.

  Differenza riassuntiva:
  \begin{itemize}
    \item \textbf{ACL}: codifica le \textit{colonne} della matrice (per ogni oggetto, chi vi accede).
    \item \textbf{Capability}: codifica le \textit{righe} della matrice (per ogni soggetto, a cosa accede).
  \end{itemize}
}

\subsection{Requisiti del meccanismo di capability}
Perche' il meccanismo di capability funzioni, si deve garantire che:
\begin{itemize}
  \item i processi \textbf{non possano forgiare} capability false;
  \item il reference monitor sia in grado di \textbf{riconoscere} se una capability e' falsa o autentica;
  \item i processi possano essere autorizzati o meno a \textbf{copiare o trasferire} le proprie capability.
\end{itemize}

\subsection{Implementazione con crittografia a chiave pubblica}
Le capability possono essere implementate tramite \textbf{crittografia a chiave pubblica}. I processi ricevono capability nella forma di triple:
\[
  \langle \texttt{object},\; \texttt{access\_rights\_for\_object},\; \texttt{unique\_code} \rangle
\]
firmate con la \textbf{chiave privata dell'oggetto}. I processi possono memorizzare e osservare le capability ma \textbf{non modificarle}, poiche' non posseggono la chiave privata per ri-firmarle (analogia con i certificati X.509).

Quando un processo vuole accedere a una risorsa, presenta la capability che detiene per quell'oggetto. L'oggetto (reference monitor) verifica:
\begin{enumerate}
  \item la \textbf{firma} digitale;
  \item il \textbf{nome} dell'oggetto;
  \item il \textbf{codice di controllo} (unique\_code);
  \item che l'accesso richiesto sia \textbf{permesso} dai diritti elencati nella capability.
\end{enumerate}
La capability puo' essere \textit{copiata} e \textit{trasferita} a un altro processo, ma \textbf{non modificata}.

\subsection{Revoca dei diritti di accesso}
La revoca puo' essere classificata lungo quattro dimensioni:
\begin{itemize}
  \item \textbf{immediata} o \textbf{ritardata};
  \item \textbf{selettiva} (solo alcuni soggetti) o \textbf{generale} (tutti);
  \item \textbf{parziale} (alcuni diritti) o \textbf{totale} (tutti i diritti);
  \item \textbf{temporanea} o \textbf{permanente}.
\end{itemize}

\textbf{Revoca in sistemi ACL}: \textit{facile} — e' sufficiente aggiornare i diritti nella lista associata all'oggetto.

\textbf{Revoca in sistemi capability}: \textit{difficile} — i diritti non sono all'oggetto ma distribuiti ai processi come capability; localizzarli puo' essere difficile o impossibile. Due approcci:
\begin{itemize}
  \item \textbf{Capability a scadenza} (\textit{time-limited}): ogni capability ha una "data di scadenza" dopo la quale va rinnovata. Non rinnovandola si ottiene una revoca (ritardata).
  \item \textbf{Capability indirette}: le capability non puntano direttamente agli oggetti ma a voci di una tabella intermedia che punta agli oggetti. Modificando la voce nella tabella si ottiene una revoca (immediata).
\end{itemize}

\section{Hybrid access control}
I sistemi reali non sono \textit{puramente} ACL-based ne' \textit{puramente} capability-based, ma \textbf{ibridi}:

\begin{center}
  \footnotesize
  \begin{tabular}{|l|l|}
    \hline
    \textbf{Funzionalita'} & \textbf{Meccanismo} \\ \hline
    Controllo accesso basato sull'identita' & ACL \\ \hline
    Facilita' di revoca & ACL \\ \hline
    Efficienza di accesso & Capability \\ \hline
  \end{tabular}
\end{center}

\subsection{UNIX come esempio di sistema ibrido}
La system call \texttt{open()} implementa il lato \textbf{ACL}:
\begin{verbatim}
int open(const char *pathname, int flags);
\end{verbatim}
dove \texttt{flags} e' \texttt{O\_RDONLY}, \texttt{O\_WRONLY} o \texttt{O\_RDWR}. Al momento della chiamata, il kernel verifica che il file esista e che l'accesso richiesto sia permesso per l'\textbf{effective-user-id} e l'\textbf{effective-group-id} del processo. Se il controllo ha successo, ritorna un intero chiamato \textbf{file descriptor}.

Il file descriptor e' un indice nella \textbf{File Descriptor Table} mantenuta dal kernel. La File Descriptor Table e' una lista di \textbf{capability}: il processo puo' usarle (puntarvi) ma non modificarle. Le successive chiamate \texttt{read()} e \texttt{write()} usano il file descriptor \textbf{senza ulteriori controlli ACL}. In questo modo il costo di verifica dell'accesso (che richiede letture su disco) e' pagato \textbf{una volta sola} e ammortizzato su migliaia di operazioni di I/O.

\section{Modelli di flusso delle informazioni}

\subsection{Bell-LaPadula (BLP)}
\dfn{Modello Bell-LaPadula}{
  Modello formale per il controllo degli accessi sviluppato negli anni '70 (Bell e LaPadula, 1973--1976) per le esigenze del Dipartimento della Difesa USA. Si concentra sulla \textbf{confidenzialita'} dei dati e sull'accesso a informazioni classificate (a differenza del modello BIBA, che descrive regole per la protezione dell'integrita').

  Ad ogni soggetto e ad ogni oggetto viene assegnata una \textbf{classe di sicurezza} (livello). Nella formulazione base le classi formano una gerarchia rigida. Esempio — schema militare statunitense (dal piu' alto al piu' basso):
  \begin{enumerate}
    \item top secret
    \item secret
    \item confidential
    \item restricted
    \item unclassified
  \end{enumerate}
}

Le regole di flusso sono:
\begin{itemize}
  \item \textbf{No read up}: un soggetto puo' leggere solo oggetti a livelli di sicurezza \textbf{minori o uguali} al suo.
  \item \textbf{No write down}: un soggetto puo' scrivere solo oggetti a livelli di sicurezza \textbf{pari o maggiori} al suo.
\end{itemize}
Il flusso delle informazioni e' \textbf{dal basso verso l'alto} e non viceversa. Acronimo: \textbf{WURD} (Write-Up Read-Down).

\ex{Bell-LaPadula}{
  Tre livelli: Top-Secret, Secret, Public. Alice e' a livello Secret. Alice \textbf{puo'} leggere documenti Secret e Public (read-down), e scrivere documenti Secret e Top-Secret (write-up). Alice \textbf{non puo'} scrivere documenti pubblici (write-down) ne' leggere documenti top-secret (read-up).
}

\subsection{Modello Biba}
\dfn{Modello Biba}{
  Modello \textbf{duale} a Bell-LaPadula, usato per la protezione dell'\textbf{integrita'} dei dati (non della confidenzialita'). Le regole sono invertite:
  \begin{itemize}
    \item Un soggetto puo' \textbf{scrivere} oggetti a livelli di sicurezza \textbf{minori o uguali} al suo.
    \item Un soggetto puo' \textbf{leggere} oggetti a livelli di sicurezza \textbf{pari o maggiori} al suo.
  \end{itemize}
  Flusso: \textbf{dall'alto verso il basso}. Acronimo: \textbf{WDRU} (Write-Down Read-Up). Nel mondo Windows prende il nome di \textit{Integrity Level} (IL).
}

\ex{Modello Biba}{
  Tre livelli: Capitano, Tenente, Carabiniere Scelto. Alice e' a livello Tenente. Alice \textbf{puo'} leggere documenti rilasciati da Tenenti e Capitani (read-up), e scrivere documenti per Tenenti e Carabinieri Scelti (write-down). In questo modo Alice non puo' dare ordini a un suo superiore.
}

\section{Permessi Linux}

\subsection{chmod}
E' possibile modificare i permessi associati a un file tramite il comando \texttt{chmod} (\textit{change file modes or Access Control Lists}):
\begin{center}
  \texttt{chmod [who] <+/-> <permissions> <file1> [file2] \dots}
\end{center}
dove \texttt{who} puo' essere \texttt{u} (user), \texttt{g} (group), \texttt{o} (other), \texttt{a} (all); il simbolo \texttt{+}/\texttt{-} aggiunge/rimuove il permesso; i permessi sono \texttt{r} (read), \texttt{w} (write), \texttt{x} (execute).

\texttt{chmod} supporta anche la notazione \textbf{ottale}: ogni permesso e' un bit:
\begin{itemize}
  \item bit 0 (valore \texttt{1}): \textbf{execute} — possibilita' di eseguire il file.
  \item bit 1 (valore \texttt{2}): \textbf{write} — possibilita' di scrivere il file.
  \item bit 2 (valore \texttt{4}): \textbf{read} — possibilita' di leggere il file.
\end{itemize}
I tre gruppi (user, group, other) danno tre cifre ottali. Es.\ \texttt{chmod 755} = \texttt{rwxr-xr-x}.

\nt{
  Impostare \texttt{777} significa regalare il file a tutti (other: read, write, execute). \textbf{Non farlo su file sensibili.}
}

\subsubsection{Permessi sulle directory}
Sulle \textbf{cartelle} i permessi hanno un significato leggermente diverso:
\begin{itemize}
  \item \textbf{execute} (bit 0): possibilita' di \textbf{entrare} in una cartella (\texttt{cd}).
  \item \textbf{write} (bit 1): possibilita' di \textbf{eliminare i file} contenuti in una cartella --- anche quelli non propri su cui non si ha permesso di scrittura!
  \item \textbf{read} (bit 2): possibilita' di \textbf{leggere il contenuto} (elenco dei file) di una cartella.
\end{itemize}
\ex{Security by obscurity con le directory}{
  \texttt{chmod 711 test/} assegna alla cartella \texttt{test/}: execute per tutti (si puo' entrare e accedere ai file se si conosce il nome), ma \textbf{niente read} per group/other (non si puo' fare \texttt{ls}). Chi non conosce il nome dei file non puo' scoprirli. E' un caso di \textit{security by obscurity} — utile come strato aggiuntivo, non come unica difesa.
}

\subsection{chown e chgrp}
E' possibile cambiare il proprietario (o il gruppo) di appartenenza di un file:
\begin{itemize}
  \item \texttt{chown [OPTION] OWNER[:GROUP] FILE} — cambia user e group proprietario.
  \item \texttt{chgrp [OPTION] GROUP FILE} — cambia solo il gruppo.
\end{itemize}
\nt{Solo \textbf{root} (normalmente) puo' regalare i file agli altri.}

\subsection{Gruppi}
Un utente puo' appartenere a \textbf{piu' gruppi}. I gruppi vengono usati (ma non solo) per controllare accessi a file di device driver (\texttt{/dev}), audio, video, stampante, dischi esterni, ecc.

\section{Permessi speciali}
Esistono 3 permessi speciali, salvati nei bit piu' significativi dei metadati del file:

\subsection{Real, effective e saved user ID}
Ogni processo ha \textbf{tre} coppie di identificatori:

\dfn{Real user/group ID}{
  \textbf{real-user-id} e \textbf{real-group-id}: identificano l'utente e il gruppo che hanno lanciato il processo. Vengono letti dal file \texttt{/etc/passwd} e \textbf{non cambiano} durante l'esecuzione del processo. Ogni processo creato eredita il real-user-id e real-group-id dell'utente. Quando un processo crea un nuovo file, proprietario e gruppo sono impostati al real-user-id e real-group-id del processo.
}

\dfn{Effective user/group ID}{
  \textbf{effective-user-id} e \textbf{effective-group-id}: sono quelli usati per \textbf{determinare i diritti di accesso} quando il processo interagisce con il filesystem. Normalmente coincidono con i real ID, ma possono cambiare dinamicamente tramite il meccanismo \textbf{setuid}.
}

\dfn{Saved user/group ID}{
  \textbf{saved-user-id} e \textbf{saved-group-id}: contengono copie degli effective ID al momento in cui un programma setuid viene eseguito. Permettono al processo di \textbf{tornare} al proprio effective user/group id precedente una volta terminata l'esecuzione del programma setuid.
}

Normalmente: effective-user-id $ = $ real-user-id e effective-group-id $ = $ real-group-id.

\subsection{Il problema di \texttt{/etc/passwd} e la soluzione setuid}
Un utente dovrebbe poter cambiare \textbf{la propria} password, ma non vedere o modificare quelle altrui. In UNIX pero' i permessi sono alla granularita' dell'\textbf{intero file}, non dei singoli record (righe di \texttt{/etc/passwd}). Non e' quindi possibile definire permessi per singola riga. L'unica alternativa senza setuid sarebbe dare a tutti read/write su \texttt{/etc/passwd} — chiunque potrebbe leggere o modificare la password di chiunque altro. La soluzione e' il meccanismo \textbf{setuid}.

\subsection{setuid}
Un file con il bit \textbf{setuid} viene eseguito con i \textbf{permessi del proprietario} del file, anziche' come utente principale. Meccanismo preciso al momento dell'esecuzione di un file con setuid:
\begin{enumerate}
  \item \textit{saved-user-id} $ \leftarrow $ \textit{effective-user-id} (salvo il vecchio effective)
  \item \textit{effective-user-id} $ \leftarrow $ user-id del proprietario del file
\end{enumerate}
I nuovi permessi restano in vigore \textbf{solo durante l'esecuzione}; al termine, il processo torna al saved-user-id. Questo permette a un processo di cambiare dominio di protezione \textbf{dinamicamente}, implementando il principio di least privilege.

\ex{setuid su \texttt{passwd}}{
  \texttt{/usr/bin/passwd} ha i permessi \texttt{-rwsr-xr-x} con proprietario \texttt{root} e \texttt{/etc/passwd} ha permessi \texttt{rw-------} (solo root). Quando un utente qualsiasi esegue \texttt{passwd}: l'\textit{effective-user-id} del processo diventa \texttt{root}, quindi il processo \textbf{puo'} scrivere \texttt{/etc/passwd} — ma solo dopo aver eseguito tutti i controlli implementati da \texttt{/bin/passwd} (verifica identita', validazione nuova password, ecc.). L'utente non ottiene accesso diretto al file.
}

\subsection{setgid}
Un file \textbf{setgid} eseguito da un utente gira con il \textbf{gruppo di appartenenza} del file. Su una cartella con setgid, i file creati al suo interno ereditano il gruppo della cartella (utile per mantenere il gruppo corretto per tutti i file in una cartella condivisa).

\subsection{Sticky bit}
Lo \textbf{sticky bit} su una cartella permette l'eliminazione dei file al suo interno solo al loro \textbf{owner}, anche se la cartella ha permesso di scrittura per tutti (es.\ \texttt{/tmp/}). Lo sticky bit su un \textbf{file} e' ignorato e deprecato.

\nt{
  Se un eseguibile con \textbf{setuid} o \textbf{setgid} appartiene a root, un altro utente potrebbe abusarne per ottenere l'accesso all'utente root. Molti eseguibili con questo bit impostati possono essere sfruttati per ottenere l'accesso a root. \textbf{ESTREMAMENTE PERICOLOSO!!!}
}

\section{Attacchi ai permessi speciali}

\subsection{PATH hijacking}
Consideriamo un file setuid che internamente chiama \texttt{system("whoami")}. Il sistema eseguira' il comando \texttt{whoami} con i privilegi dell'utente indicato. \texttt{whoami} e' un comando shell, che puo' essere \textbf{dirottato} modificando il \texttt{PATH}:
\begin{verbatim}
PATH=. ./vulnprogram
\end{verbatim}
Se nella directory corrente esiste un file chiamato \texttt{whoami} con contenuto arbitrario, il sistema lo eseguira' \textit{al posto} di \texttt{/usr/bin/whoami}, con i privilegi elevati del programma setuid.

\subsection{sudo e privilege elevation}
Il comando \texttt{sudo} esegue comandi con privilegi elevati e controlla la password dell'utente tramite un framework chiamato \textbf{PAM}.

\ex{CVE-2023-22809 — sudoedit}{
  In Sudo versioni 1.8.0--1.9.12p1, la funzionalita' \texttt{sudoedit} (\texttt{-e}) gestisce in modo scorretto gli argomenti extra passati nelle variabili d'ambiente \texttt{SUDO\_EDITOR}, \texttt{VISUAL} ed \texttt{EDITOR}. Cio' permette a un attaccante locale di aggiungere voci arbitrarie alla lista dei file da processare, portando a privilege escalation. Il problema esiste perche' un editor specificato dall'utente puo' contenere un argomento \texttt{--} che vanifica un meccanismo di protezione (es.\ \texttt{EDITOR="vim -- /path/to/extra/file"}). Severity CVSS 3.x: \textbf{7.8 HIGH}.
}

% \end{document}
